
\documentclass{beamer}

\usepackage{xcolor}
\usepackage{color}
\usepackage{graphicx}
\def\xcolorversion{2.00}
\def\xkeyvalversion{1.8}
\setbeamersize{text margin left=0pt,text margin right=0pt}

\makeatletter
\newcommand*\bigcdot{\mathpalette\bigcdot@{.9}}
\newcommand*\bigcdot@[2]{\mathbin{\vcenter{\hbox{\scalebox{#2}{$\m@th#1\bullet$}}}}}
\makeatother

\usepackage{TIKZ_Macra}
\usepackage{xparse}
\usepackage[utf8]{inputenc}
\usepackage[T1]{fontenc}

\useinnertheme[shadow=true]{rounded}
\setbeamercolor{block title}{use=structure,fg=structure.fg,bg=structure.fg!30!bg}
\setbeamercolor{block body}{parent=normal text,use=block title,bg=block title.bg!50!bg}
\setbeamercolor{block title example}{use=example text,fg=example text.fg,bg=example text.fg!30!bg}
\setbeamercolor{block body example}{parent=normal text,use=block title example,bg=block title example.bg!50!bg}

\usepackage[most]{tcolorbox}
\usepackage{lmodern}
\usepackage{mathtools}
\usepackage{amsfonts}
\usepackage{amssymb}
\usepackage{latexsym}
\usepackage{hyperref}
\usepackage{changepage}

\newtcbox{\mybox}{blank, on line, opacitytext=0.5}

\usepackage{Moje_makra}
\usepackage{Moje_Zmienne}
\usepackage{Zbiory_i_konfiguracje}
\usepackage{Complexity}
\usepackage{Sieci}

\newcommand{\Tr}{\mathsf{Tr}}
\newcommand{\Runs}{\mathsf{Runs}}
\newcommand{\PreA}{\mathsf{Pre}}
\newcommand{\PostA}{\mathsf{Post}}
\newcommand{\Part}{\mathcal{P}}
\newcommand{\Blocks}{\mathsf{Blocks}}
\newcommand{\Sat}{\mathsf{Sat}}
\newcommand{\true}{\mathsf{true}}
\newcommand{\false}{\mathsf{false}}
\newcommand{\Act}{\Sigma}
\newcommand{\AP}{\AProp}
\newcommand{\impl}{\mathsf{Impl}}
\newcommand{\absys}{\mathsf{Abs}}
\newcommand{\depth}{\mathsf{md}}

\begin{document}

\begin{frame}[fragile]{}
\Margines{
\begin{center}
  {\huge Theory of Concurrency}\medskip

  {\large Lecture 3a: Recap of temporal logics}\medskip

  Material from the previous lecture\medskip

  Piotr Hofman
\end{center}
}
\end{frame}


\begin{frame}{What was covered}
\Margines{
\begin{enumerate}
  \item Kripke structures as state-based models.
  \item LTL as a logic of linear infinite behaviours.
  \item Translation from LTL to \Buchi\ automata.
  \item LTL model checking by emptiness of a product automaton.
  \item CTL and CTL* as branching-time logics.
\end{enumerate}

\begin{block}{Key distinction}
LTL observes complete traces. CTL and CTL* observe the branching structure of possible futures.
\end{block}
}
\end{frame}

\begin{frame}{Kripke structures and traces}
\Margines{
\begin{block}{Kripke structure}
A Kripke structure over atomic propositions \(\AP\) is
\[
  M=(S,I,R,L),
\]
where \(S\) is a finite set of states, \(I\subseteq S\) is the set of initial states, \(R\subseteq S\times S\) is left-total, and
\[
  L:S\to 2^{\AP}.
\]
\end{block}

\begin{block}{Trace of a path}
For a path \(\pi=s_0s_1s_2\ldots\), its trace is
\[
  L(s_0)L(s_1)L(s_2)\ldots\in (2^{\AP})^\omega.
\]
Thus, for LTL over \(M\), the alphabet is
\[
  \Sigma=2^{\AP}.
\]
\end{block}
}
\end{frame}

\begin{frame}{LTL syntax}
\Margines{
\begin{block}{Definition}
An LTL formula is generated by
\[
\varphi ::= \true \mid p \mid \neg\varphi \mid
\varphi_1\land\varphi_2 \mid X\varphi \mid \varphi_1 U \varphi_2,
\]
where \(p\in\AP\).
\end{block}

\begin{block}{Derived operators}
\[
F\varphi \equiv \true U\varphi,
\qquad
G\varphi \equiv \neg F\neg\varphi.
\]
\begin{itemize}
  \item \(F\varphi\): eventually \(\varphi\).
  \item \(G\varphi\): always \(\varphi\).
\end{itemize}
\end{block}
}
\end{frame}

\begin{frame}{LTL semantics}
\Margines{
Let \(w=w_0w_1w_2\ldots\in(2^{\AP})^\omega\). We write \(w,i\models\varphi\).

\begin{block}{Temporal clauses}
\begin{align*}
 w,i &\models X\varphi
 &&\text{iff } w,i+1\models\varphi,\\
 w,i &\models \varphi U\psi
 &&\text{iff there is } j\geq i \text{ such that } w,j\models\psi \\
 &&&\text{and }w,k\models\varphi\text{ for all }i\leq k<j.
\end{align*}
\end{block}

\begin{block}{Model satisfaction}
\[
M,s\models_{LTL}\varphi
\quad\text{iff every trace starting from }s\text{ satisfies }\varphi.
\]
\end{block}
}
\end{frame}

\begin{frame}{LTL as a language specification}
\Margines{
Each LTL formula defines an \(\omega\)-language:
\[
  \words{\varphi}=\{w\in(2^{\AP})^\omega : w\models\varphi\}.
\]

\begin{block}{Verification condition}
For a state \(s\), let \(\Tr(s)\) be the set of all infinite traces from \(s\). Then
\[
  M,s\models_{LTL}\varphi
  \quad\Longleftrightarrow\quad
  \Tr(s)\subseteq \words{\varphi}.
\]
\end{block}

\begin{block}{Consequence}
LTL is insensitive to branching that does not change the set of traces.
\end{block}
}
\end{frame}

\begin{frame}{LTL to \Buchi\ automata}
\Margines{
\begin{block}{Theorem}
For every LTL formula \(\varphi\), one can construct a \Buchi\ automaton \(B_\varphi\) of exponential size such that
\[
  \lang{B_\varphi}=\words{\varphi}.
\]
\end{block}

\begin{block}{Model checking idea}
To check \(M,s\models\varphi\), build an automaton \(A_s\) for the traces of \(M\) from \(s\), then check
\[
  \lang{A_s}\cap \lang{B_{\neg\varphi}}=\emptyset.
\]
\end{block}

\begin{block}{Interpretation}
A non-empty product gives a counterexample trace.
\end{block}
}
\end{frame}

\begin{frame}{CTL and CTL* in one slide}
\Margines{
\begin{block}{CTL}
CTL alternates path quantifiers with temporal operators:
\[
EX\varphi,\quad AX\varphi,\quad E(\varphi U\psi),\quad A(\varphi U\psi),\quad EF\varphi,\quad AG\varphi.
\]
\end{block}

\begin{block}{CTL*}
CTL* separates state formulas and path formulas. It allows arbitrary LTL combinations under path quantifiers:
\[
A\alpha,
\qquad
E\alpha.
\]
\end{block}

\begin{block}{Important point}
CTL and CTL* observe not only which traces exist, but also how they branch from the current state.
\end{block}
}
\end{frame}

\begin{frame}{Why this matters for Lecture 3}
\Margines{
\begin{block}{For LTL}
A state can be abstracted by its trace language:
\[
  \Tr(s)\subseteq(2^{\AP})^\omega.
\]
Trace equivalence preserves all LTL formulas.
\end{block}

\begin{block}{For CTL and CTL*}
The trace language is not enough, because the formulas may ask about possible choices from the current state.
\end{block}

\begin{block}{Next question}
What is the right abstraction for branching-time logics?
\end{block}
}
\end{frame}

\end{document}
